%%%%%%%%%%%%%%%%%%%%%%%
%%% DOCUMENT SETTINGS
%%%%%%%%%%%%%%%%%%%%%%%
\documentclass[12pt]{exam}
\printanswers % Uncomment to show answers
\include{formatAndDefs}
%%%%%%%%%%%%%%
%%% PACKAGES
%%%%%%%%%%%%%%
\usepackage{amsmath,amsfonts,amssymb,amsthm} % math fonts and symbols
\usepackage{bbm} % Math blackboard font
\usepackage{enumitem} % For reference enumerations lists
\usepackage{textcomp} % some symbols (like copyright)
\usepackage[top = 2cm, bottom = 2cm, right = 1.2cm, left = 1.2cm, paper = a4paper, headheight = 6cm]{geometry} % Margins setup
\usepackage{xcolor} % Colors
\usepackage{pgfplots}
\usepackage{spalign}
%%%%%%%%%%%%%%%%%%%%%%%%%%
%%% MACROS AND COMMANDS
%%%%%%%%%%%%%%%%%%%%%%%%%%
\definecolor{CUNEForange}{RGB}{237, 114, 1}
\newcommand{\N}{\mathbb{N}} % natural number set
\newcommand{\R}{\mathbb{R}} % real number set
%%%%%%%%%%%%%%%%%%%%%%%%%%%%%%%%%
%%% HEADER AND FOOT 
%%%%%%%%%%%%%%%%%%%%%%%%%%%%%%%%%
% Defining information
\newcommand{\degree}[1]{Bachelor in Philosophy, Politics and Economics}
\newcommand{\shortdeg}[1]{PPE}
\newcommand{\exam}[1]{Final Exam}
\newcommand{\subject}[1]{Mathematics}
\newcommand{\theyear}[1]{2023/24}
%\newcommand{\ccauthor}[1]{A. Azze \textcopyright\the\year}
\newcommand{\thedate}[1]{10 06 2024}
\newcommand{\university}[1]{CUNEF University}
% Setting header
\usepackage{fancyhdr}
\pagestyle{fancy}
\renewcommand{\headrulewidth}{0.1mm}
\renewcommand{\footrulewidth}{0.1mm}
% \renewcommand{\sectionmark}[1] % To avoid sections' titles being displayed on the header
{\markright{\MakeUppercase{}}}
\setlength\headsep{0.5cm}
% Header and footer
\fancyhead[L]{\degree{} \\ \subject{} - \theyear{}}
\fancyhead[R]{\thedate{} \\ \ccauthor{}}
\fancyfoot[L]{\university{}}
\fancyfoot[C]{}
\fancyfoot[R]{\thepage}

%%%%%%%%%%%%%%%%%%%%%%%%%
%%% DOCUMENT CONTENT
%%%%%%%%%%%%%%%%%%%%%%%%%
\begin{document}
%--- Instructions ---%
\begin{center}

    {\Huge \textbf{\exam{}}}
    \vspace{0.7cm}

    \begin{flushright}
        \textbf{Time Limit:} 120 minutes
    \end{flushright}
    \vspace{-0.2cm}
    
    \fbox{\parbox{\textwidth}{\centering 
        \begin{enumerate}[leftmargin = 0.5cm, label = $\bullet$, rightmargin = 0.2cm]
    
            \item \textbf{Every non-trivial calculation and claim in the exam must be properly justified.}

            \item The use of calculators is not necessary and therefore not allowed.

            \item Digital devices such as mobile phones, tablets, and laptops, as well as internet access, are forbidden.

            \item You can use physical resources such as books, notes, and other printed materials to assist you during the exam.
 
        \end{enumerate}
    }}
\end{center}
%--- Instructions ---%

%--- Questions ---%
\noindent
\begin{questions}

\question[1] Given the following supply and demand sets,
\[S=\left\lbrace 2p+q=5 \right\rbrace\qquad D=\left\lbrace 3p-q=8\right\rbrace.\]
Plot the demand and supply sets and compute the equilibrium point. Explain its meaning using the plot.
\question[2] Compute the following limits.
\begin{itemize}
    \item[i)] $\displaystyle\lim_{x\rightarrow 1}\frac{x^2-1}{x+1}$
    \item[ii)] $\displaystyle\lim_{x\rightarrow\infty}\frac{x^2-x+1}{2x^2+3x-1}$
\end{itemize}
\question[1] Compute the global maximum and minimum, if they exist, of the function $$f(x)=x^3-\frac{3}{2}x^2$$ in the interval $[-2,2]$. How do they change if we expand the interval to $[-10,10]$?

\question[1] Determine the maxima and minima, if they exist, of the function $$f(x,y)=3x^2-2xy-y^2-8y.$$

\question[2] Compute the following integrals.
\begin{itemize}
    \item[i)] $\displaystyle \int \frac{1}{(x-3)(x-2)} dx$
    \item[ii)] $\displaystyle \int_1^{\infty} 2xe^{-3x^2} dx$
\end{itemize}
\question[3] Given the following matrices.
\begin{equation*}
A=\begin{pmatrix}
1 & 0 & 0 \\
0 & 0 & 1 \\
0 & 1 & 0
\end{pmatrix},
\enspace
B=\begin{pmatrix}
1 & 2 & 0\\
1 & -2 & -1\\
2 & 0 & -1
\end{pmatrix},
\enspace
t=\begin{pmatrix}
x\\
y\\
z
\end{pmatrix},
\enspace
c=\begin{pmatrix}
0\\
1\\
1
\end{pmatrix}.
\end{equation*}
\begin{itemize}
    \item[i)] Compute $A^{1245}$
    \item[ii)] Determine the value of the matrix $C$ that satisfies $C\cdot A=B$
    \item[iii)] Write the linear system $Bt=c$ and solve it using Gaussian elimination.
\end{itemize}
\end{questions}	
\end{document}